% संपूर्ण मराठी ग्रंथातील प्रकरण 73: विधाने म्हणजे प्रकार
% हा संपादनयोग्य स्रोतखंड आहे. संकलनासाठी संपूर्ण स्रोतसंग्रहातील
% अक्षररूपे, पूर्वभाग, चिन्हव्याख्या आणि आकृती-स्रोत आवश्यक आहेत.
% मूळ: Open Logic Project; मजकूर आणि रूपांतर CC BY 4.0.
% यंत्रानुवाद, दुरुस्ती आणि तपासणी: OpenAI Codex —
% GPT-5.6 Sol आणि GPT-6 Sol, दोन्ही Ultra effort.
% दुरुस्ती-पुनरावलोकन: GPT-6.1 Sol, Ultra effort.
\chapter{विधाने म्हणजे प्रकार}\label{int:pty:chap}
\begin{editorial}
  करी--हॉवर्ड अनुरूपतेवरील या प्रकरणाचा हा \emph{अतिशय प्रायोगिक}
  मसुदा आहे. याला अधिक स्पष्टीकरण आणि प्रेरणा हवी; यात चुका
  व गाळलेले भागही असण्याची शक्यता आहे. सामान्यीकरणाची सिद्धता
  तपासून विस्तारित केली पाहिजे. गुणाकार-प्रकाराची उदाहरणे नाहीत.
  क्रमपरिवर्तन आणि सरलीकरण रूपांतरणांचा समावेश केलेला नाही.
  येथे मूलतः गृहीत धरलेल्या (प्रकारयुक्त) लॅम्डा कलनाविषयीही
  साहित्य आल्यावर हे प्रकरण अधिक समजेल. अत्यंत सावधगिरीने वापरा.
\end{editorial}









\section{परिचय}\label{int:pty:int:sec}

लॅम्डा कलन हे संगणनाचे एक साधे प्रतिमान आहे. ते चरांपासून
बनलेल्या पदांवर कार्य करते. $N$ आणि $M$ ही पदे असतील, तर
$(NM)$ हेही एक पद आहे (``$N$ चे~$M$ वर उपयोजन''). $N$ हे पद
असेल, तर $\lambd[x][N]$ हेही पद आहे. $N$ मध्ये~$x$ हे चर
असेल, तर $\lambd[x][N]$ मधील~$x$ ची संबंधित स्थाने
$\lambd[x]$ या संकारकाने बद्ध होतात; त्यामुळे ती मुक्त राहत
नाहीत. $(\lambd[x][N]M)$ या रूपातील पदाचे
$\Subst{N}{M}{x}$ मध्ये $\beta$-संकुचन होते असे म्हणतात
(म्हणजे~$N$ मधील~$x$ ची प्रत्येक मुक्त आस्थिती~$M$ ने
बदलल्यावर मिळणारे पद). $N$ मधील एखाद्या उपपदाचे
$\beta$-संकुचन करून~$N'$ मिळत असेल, तर~$N$ चे~$N'$ मध्ये
$\beta$-परिवर्तन होते, आणि आपण $N \redone N'$ असे लिहितो.
$\lambd$-कलनातील संगणन म्हणजे
$N \redone N_1 \redone N_2 \redone \dots$ असा क्रम.
हा क्रम~$N_k$ या अशा पदावर संपला की त्याच्या कोणत्याही
उपपदाचे $\beta$-संकुचन होऊ शकत नाही, तर त्या पदाला
\emph{सामान्य} किंवा~$N$ चे सामान्य रूप म्हणतात.
$\lambd$-कलनातील संगणनाचा विचार अशा क्रमांच्या रूपात करा.
सामान्य रूप मिळाल्यास संगणन थांबते, आणि ते सामान्य रूपच
संगणनाचे फल असते. हे साधे संगणन-प्रतिमान ट्यूरिंग-संपूर्ण
ठरते: प्रत्येक आंशिक पुनरावर्ती फलनाची गणना
$\lambd$-कलनातील एका पदाने करता येते.

हास्केल करी आणि विल्यम हॉवर्ड यांनी स्वतंत्रपणे नैसर्गिक
निगमन आणि $\lambd$-कलन यांतील जवळचे साम्य शोधले.
सहज अर्थाने $\lambd[x][N]$ हे पद एक फलन आहे: $x$ हा त्याचा
आर्ग्युमेंट, आणि~$x$ दिल्यावर मूल्य कसे मोजायचे हे~$N$
सांगते. म्हणून $(\lambd[x][N]M)$ म्हणजे
$\lambd[x][N]$ हे फलन~$M$ या आर्ग्युमेंटवर लावणे;
$\beta$-संकुचनाने त्याचे मूल्यांकन कसे करायचे ते कळते:
$N$ मध्ये~$x$ ऐवजी~$M$ ठेवायचे (आणि मिळालेल्या पदाचे
मूल्यांकन पुढे चालू ठेवायचे).
आता या फलनांना निश्चित प्रांत आणि मूल्यसंच आहेत असे समजा.
$\lambd[x][N]$ चा प्रांत~$A$ आणि मूल्यसंच~$B$ असेल, तर
$x$ ला फक्त~$A$ मधील आर्ग्युमेंट घेता येतील आणि~$N$ ने
$B$ मधील मूल्ये द्यावीत. त्यामुळे $\lambd[x][N]$ हे
$A \to B$ असे फलन आणि $M \in A$ असेल, तेव्हाच
$(\lambd[x][N]M)$ ला अर्थ असावा. ही माहिती
$\lambd$-कलनात घातली की \emph{प्रकारयुक्त} $\lambd$-कलन
मिळते. प्रकारयुक्त $\lambd$-कलनात प्रत्येक
$(\lambd[x][N]M)$ अभिव्यक्ती ग्राह्य नसते;
$\lambd[x][N]$ आणि~$M$ यांचे ``प्रकार'' जुळतात,
म्हणजे $\lambd[x][N]$ चा प्रकार $A \lif B$ आणि~$M$ चा
प्रकार~$A$ असतो, तेव्हाच ती ग्राह्य असते.
$x$ चा प्रकार~$A$ आणि~$N$ चा प्रकार~$B$ असेल,
तेव्हाच $\lambd[x][N]$ चा प्रकार $A \to B$ असतो.
सर्वसाधारणपणे प्रत्येक $(M_1M_2)$ अभिव्यक्ती ग्राह्य नसते.
$(M_1M_2)$ हे पद ग्राह्य होण्यासाठी~$M_1$ चा प्रकार $A \to B$
आणि~$M_2$ चा प्रकार~$A$ असावा;
अशा वेळी $(M_1M_2)$ चा प्रकार~$B$ असतो.

हे प्रकार-नियम आता नैसर्गिक निगमनातील
निष्पत्ती शी स्पष्टपणे जुळतात: प्रकारांना
सूत्रे, आणि निष्पत्ती ना $\lambd$-पदे
अनुरूप असतात. उदाहरणार्थ, $A \lif B$ ची
निष्पत्ती $\lambd[\typeof{x}{A}][N]$
या पदाशी जुळते; त्याचा फलन-प्रकार $A \lif B$ आहे.
नैसर्गिक निगमनाचे अनुमान-नियम प्रकारयुक्त लॅम्डा कलनातील
प्रकार-नियमांशी जुळतात; उदाहरणार्थ,
\begin{prooftree}
  \AxiomC{$\Discharge{A}{x}$}
  \DeduceC{$B$}
  \DischargeRule{\Intro\lif}{x}
  \UnaryInfC{$A \lif B$}
  \DisplayProof\bottomAlignProof
  \quad\text{याला अनुरूप आहे}\quad
  \Axiom$x: A \fCenter N: B$
  \UnaryInf$ \fCenter \lambd[\typeof{x}{A}][N] : A \lif B$
\end{prooftree}
उजवीकडील नियमाचा अर्थ असा की~$x$ चा प्रकार~$A$ आणि
$N$ चा प्रकार~$B$ असेल, तर
$\lambd[\typeof{x}{A}][N]$ चा प्रकार~$A \lif B$ असतो.

$\Elim\lif$ हा नियम संयुक्त पदांसाठीच्या प्रकार-नियमाशी
जुळतो; म्हणजे,
\[
  \AxiomC{$A \lif B$}
  \AxiomC{$A$}
  \RightLabel{\Elim\lif}
  \BinaryInfC{$B$}
  \DisplayProof
  \quad\text{याला अनुरूप आहे}\quad
  \Axiom$\fCenter M_1: A \lif B$
  \Axiom$\fCenter M_2: A$
  \BinaryInf$ \fCenter (M_1M_2) : B$
  \DisplayProof
\]
\Intro\lif{} या नियमानंतर लगेच \Elim\lif{} नियम
आला, तर निष्पत्ती सोपी करता येते:
\begin{gather*}
  \AxiomC{$\Discharge{A}{x}$}
  \DeduceC{$B$}
  \DischargeRule{\Intro\lif}{x}
  \UnaryInfC{$A \lif B$}
  \AxiomC{}
  \DeduceC{$A$}
  \RightLabel{\Elim\lif}
  \BinaryInfC{$B$}
  \DisplayProof
  \quad\redone\quad
  \AxiomC{}
  \DeduceC{$A$}
  \DeduceC{$B$}
  \DisplayProof
\end{gather*}
हे लॅम्डा पदांच्या $\beta$-संकुचनाशी जुळते:
\[
(\lambd[\typeof{x}{A}][N]M) \quad\redone\quad
\Subst{N}{M}{x}.
\]

अशीच अनुरूपता $\land$ च्या नियमांचा ``गुणाकार''
प्रकारांशी आणि $\lor$ च्या नियमांचा ``बेरीज'' प्रकारांशीही आहे.

साध्या प्रकारयुक्त लॅम्डा कलनातील पदे आणि नैसर्गिक निगमनातील
निष्पत्ती यांतील या अनुरूपतेला ``करी--हॉवर्ड'',
``सिद्धता म्हणजे कार्यक्रम'', किंवा ``विधाने म्हणजे प्रकार''
अनुरूपता म्हणतात. सूत्रे (विधाने) प्रकारांशी आणि
सिद्धता पदांशी जुळतात; त्याशिवाय इतर अनुरूपता पुढीलप्रमाणे:
\begin{center}
  \begin{tabular}{c c c}
    तर्कशास्त्र & कार्यक्रम \\
    \hline
    विधान/सूत्र & प्रकार \\
    सिद्धता & पद \\
    गृहीतक & चर \\
    निरसित गृहीतक & बद्ध चर \\
    अनिरसित गृहीतक & मुक्त चर \\
    $A \lif B$ & फलन-प्रकार \\
    $A \land B$ & गुणाकार-प्रकार \\
    $A \lor B$ & बेरीज-प्रकार \\
    $\lfalse$ & रिक्त प्रकार \\
    \hline
  \end{tabular}
\end{center}

करी--हॉवर्ड अनुरूपता ही स्वयंचलित सिद्धता-साहाय्यकांची आणि
कार्यक्रमांसाठीच्या प्रकार-तपासकांची एक आधारभूत कल्पना आहे.
कारण एखादे विधान सिद्ध करणारी सिद्धता तपासणे (जसे आपण वर केले)
म्हणजे एखाद्या कार्यक्रमाला (पदाला) जाहीर केलेला प्रकार आहे का
हे तपासणे होय.












\section{सिद्धता-पदे}\label{int:pty:ter:sec}

\Log{N1i} मध्ये गृहीतकांना निरसन-खुणा म्हणून चलांची
चिन्हे देतो (उदा., $A^x$). \Log{N2i} मध्ये प्रत्येक
$\Gamma \Sequent A$ क्रमवर्तीच्या डावीकडील संदर्भ~$\Gamma$
हा सुसंगत चल–सूत्र जोड्यांचा संच असतो. त्यामुळे
\Log{N2i} मधील सिद्धता च्या प्रत्येक क्रमवर्तीवर
आतापर्यंतच्या सिद्धता ने काय सिद्ध केले आहे
(म्हणजे~$A$) एवढीच नव्हे, तर त्या टप्प्यावर कोणती
गृहीतके खुली आहेत (म्हणजे~$\Gamma$) ही माहितीही असते.
$A$ \emph{कसे} सिद्ध झाले याची माहितीही प्रत्येक
क्रमवर्तीत घातली तर? त्यासाठी \emph{पद-चिन्हांकित}
\Log{tN2i} पद्धती परिभाषित करता येते; तिच्या
क्रमवर्तींचे रूप $\Gamma \Sequent N : A$ असते.
$\Gamma$ आणि~$A$ यांचे अर्थ आधीप्रमाणेच आहेत:
$x: B$ या रूपातील जोड्यांचा संच, आणि क्रमवर्तीवर
संपणाऱ्या उप-सिद्धता ने~$\Gamma$ मधील गृहीतकांवरून
निष्पन्न होते असे दाखवलेले सूत्र. ~$N$ हे त्या
क्रमवर्तीवर संपणाऱ्या सिद्धता चे सांकेतिकीकरण करणारे
विशिष्ट पद असेल.

प्रत्येक क्रमवर्तीला देण्यात येणारी पदे $x$, $y$,
\dots या चलांपासून आणि आतापर्यंत झालेल्या अनुमानांचे
सांकेतिकीकरण करणाऱ्या फलन-चिन्हांपासून बनतात.
प्रत्येक अनुमान-नियमासाठी स्वतंत्र फलन-चिन्ह लागते.
परिचय नियमांना अनुरूप फलन-चिन्हांना
``रचनाकार'' आणि निरसन नियमांना अनुरूप चिन्हांना
``विघटक'' म्हणतात. नियमाची जितकी आधारविधाने, तितके
प्रत्येक फलन-चिन्हाचे आर्ग्युमेंट असतात. उदाहरणार्थ,
\Intro{\land} आणि \Elim{\land}[1] साठी अनुक्रमे
$c^\land$ (दोन-स्थानी) आणि $d^\land_1$ (एक-स्थानी)
ही फलन-चिन्हे वापरता येतील. \Log{tN2i} मध्ये हे
नियम असे होतील:
\[
\Axiom$\Gamma_1 \fCenter N: A$
\Axiom$\Gamma_2 \fCenter M: B$
\RightLabel{\Intro{\land}}
\BinaryInf$\Gamma_1, \Gamma_2 \fCenter c^\land(N, M) : A \land B $
\DisplayProof
\quad
\Axiom$\Gamma \fCenter N: A \land B$
\RightLabel{\Elim{\land}[1]}
\UnaryInf$\Gamma \fCenter d^\land_1(N) : A$
\DisplayProof
\]
अनुमान-नियम एखादे गृहीतक निरसित करत असेल,
म्हणजे आधारविधानात असलेले खुणित सूत्र~$x:A$
निष्कर्षाच्या संदर्भात नसेल, तर ते सिद्धता-पदातही
दर्शवावे लागते. अशा नियमांशी संबंधित फलन-चिन्हे
त्या चलांना \emph{बद्ध} करतात. उदाहरणार्थ,
\Intro{\lif} नियम $x:A, \Gamma \Sequent B$ या
आधारविधानापासून $\Gamma \Sequent A \lif B$
या निष्कर्षापर्यंत जातो. रचनाकार~$c^{\lif}$ एक
आर्ग्युमेंट घेतो, पण चल~$x$ बद्ध होतो हेही त्याने
दर्शवले पाहिजे. खूण~$x$ असलेले गृहीतक निरसित
झाले आहे हे सिद्धता-पद अशा प्रकारे व्यक्त करते.
उदाहरणार्थ, आधारविधानाचे सिद्धता-पद~$N$ असेल,
तर निष्कर्षाचे सिद्धता-पद $c^{\lif}([x]N)$ आणि
नियम असा लिहिता येईल:
\[
\Axiom$x:A, \Gamma \fCenter N: B$
\RightLabel{\Intro{\lif}}
\UnaryInf$\Gamma \fCenter c^{\lif}([x]N) : A \lif B$
\DisplayProof
\]

सिद्धता पूर्णपणे पुन्हा उभारण्यासाठी आणखी
काही माहिती लागते. नियमाच्या आधारविधानांतील
सूत्रे वरून निष्कर्षातील सूत्र पूर्णपणे
ठरत नसेल, तर ती माहिती पदात जोडली पाहिजे.
\Intro{\lif} मध्ये असे होते; म्हणून आपण
$c^{\lif}([x:A]N)$ असे लिहू. \Intro{\lor}
आणि~\FalseInt मध्येही असेच होते; म्हणून ती
माहिती वाहून नेणारी फलन-चिन्हे वापरू, उदाहरणार्थ:
\[
\Axiom$\Gamma \fCenter N:A$
\RightLabel{\Intro{\lor}[1]}
\UnaryInf$\Gamma \fCenter c^{\lor{B}}_1(N):A \lor B$
\DisplayProof
\quad
\Axiom$\Gamma \fCenter N:\lfalse$
\RightLabel{\FalseInt}
\UnaryInf$\Gamma \fCenter d^\lfalse_{A}(N):A$
\DisplayProof
\]

या कल्पनांवरून क्रमवर्ती-शैलीतील नैसर्गिक निगमनाची
अशी पद्धती बांधता येते की प्रत्येक क्रमवर्तीचे रूप
$\Gamma \Sequent N : A$ असते आणि~$N$ हे
\emph{सिद्धता-पद} असते.

अशी सिद्धता-पदे आणि तथाकथित \emph{प्रकारयुक्त लॅम्डा कलन}
यातील पदे यांचा जवळचा संबंध आहे. तो लक्षात घेऊन वरची
सर्वसाधारण रचनाकार-विघटक चिन्हे न वापरता लॅम्डा
कलनाच्या साहित्यातील चिन्हे वापरू. उदाहरणार्थ,
$c^{\lif}([x:A]N)$ ऐवजी $\lambd[\typeof{x}{A}][N]$,
$d^{\lif}(N,M)$ ऐवजी $(NM)$,
$c^\land(N,M)$ ऐवजी $\pair{N}{M}$, आणि
$d^\land_i(N)$ ऐवजी~$\proj{i}{N}$ लिहू.

\begin{defn}
  सिद्धता-पदे पुढील नियमांनी विगमनाने तयार होतात:
  \begin{enumerate}
  \item प्रत्येक चर~$x$ हे सिद्धता-पद आहे (स्वयंसिद्धके).
  \item $x$ हे चर, $A$ हे सूत्र, आणि~$N$ हे
    सिद्धता-पद असेल, तर $\lambd[\typeof{x}{A}][N]$
    हेही सिद्धता-पद आहे~(\Intro\lif).
  \item $N$ आणि~$M$ ही सिद्धता-पदे असतील, तर $(NM)$
    हेही सिद्धता-पद आहे~(\Elim\lif).
  \item $N$ आणि~$M$ ही सिद्धता-पदे असतील, तर
    $\pair{N}{M}$ हे सिद्धता-पद आहे~(\Intro\land).
  \item $N$ हे सिद्धता-पद असेल, तर $\proj{i}{N}$
    हेही सिद्धता-पद आहे; येथे~$i$ हे~$1$ किंवा~$2$
    असते~(\Elim\land[i]).
  \item $N$ हे सिद्धता-पद आणि~$A$ हे सूत्र असेल,
    तर $\inj[A]{i}{N}$ हे सिद्धता-पद आहे; येथे~$i$
    हे~$1$ किंवा~$2$ असते~(\Intro\lor[i]).
  \item $M$, $N_1$, $N_2$ ही सिद्धता-पदे आणि~$x_1$,
    $x_2$ ही चले असतील, तर
    $\dcase{M}{x_1}{N_1}{x_2}{N_2}$ हे सिद्धता-पद
    आहे~(\Elim\lor).
  \item $N$ हे सिद्धता-पद आणि~$A$ हे सूत्र
    असेल, तर $\abort{A}{N}$ हे सिद्धता-पद आहे~(\FalseInt).
  \end{enumerate}
\end{defn}

प्रत्येक सिद्धता-पद हे एखाद्या सिद्धता चे
सिद्धता-पद असेलच असे नाही. उदाहरणार्थ,
$\pair{N}{M}$ हे पद~\Intro{\land} अनुमानावर
संपणाऱ्या सिद्धतेचे सांकेतिकीकरण करते; म्हणून तिचे
अंतिम सूत्र $A \land B$ हे संयोगसूत्र असते.
हे~\Elim{\lif} नियमाचे मुख्य आधारविधान होऊ शकत
नाही; त्यामुळे $(\pair{N}{M}P)$ हे योग्य सिद्धतेचे
सिद्धता-पद होऊ शकत नाही. \Log{tN2ip} च्या नियमांचे
पालन करणारी सिद्धता-पदेच \emph{योग्य} सिद्धता-पदे
आहेत. हे नियम \ref{int:pty:ter:tab:tN2ip} मध्ये दिले आहेत.








\begin{table}
\begin{tabular}{@{}c@{}c@{}}
\multicolumn{2}{l}{\textbf{स्वयंसिद्धके:}}\\
\multicolumn{2}{c}{$x:A \Sequent x:A$}\\[3ex]
\multicolumn{2}{l}{\textbf{अनुमान-नियम:}}\\[3ex]
\multicolumn{2}{c}{\Axiom$x: A, \Gamma \fCenter N:B$
\RightLabel{\Intro{\lif}}
\UnaryInf$\Gamma \fCenter \lambd[\typeof{x}{A}][N] : A \lif B $
\DisplayProof}
\\[4ex]
\multicolumn{2}{c}{\Axiom$\Gamma_1 \fCenter N:A \lif B$
\Axiom$\Gamma_2 \fCenter M:A$
\RightLabel{\Elim{\lif}}
\BinaryInf$\Gamma_1, \Gamma_2 \fCenter (NM) : B$
\DisplayProof}
\\[4ex]
\multicolumn{2}{c}{\Axiom$\Gamma_1 \fCenter N_1 : A$
\Axiom$\Gamma_2 \fCenter N_2 : B$
\RightLabel{\Intro{\land}}
\BinaryInf$\Gamma_1, \Gamma_2 \fCenter \pair{N_1}{N_2} : A \land B $
\DisplayProof}
\\[4ex]
\Axiom$\Gamma \fCenter N: A \land B$
\RightLabel{\Elim{\land}[1]}
\UnaryInf$\Gamma \fCenter \proj{1}{N} : A$
\DisplayProof
&
\Axiom$\Gamma \fCenter N:A \land B$
\RightLabel{\Elim{\land}[2]}
\UnaryInf$\Gamma \fCenter \proj{2}{N} : B$
\DisplayProof\\[3ex]
\\[4ex]
\Axiom$\Gamma \fCenter N: A$
\RightLabel{\Intro{\lor}[1]}
\UnaryInf$\Gamma \fCenter \inj[B]{1}{N} : A \lor B$
\DisplayProof
&
\Axiom$ \Gamma \fCenter N : B$
\RightLabel{\Intro{\lor}[2]}
\UnaryInf$\Gamma \fCenter \inj[A]{2}{N} : A \lor B$
\DisplayProof\\[3ex]
\multicolumn{2}{c}{
\begin{tabular}[b]{@{}l@{}}
\Axiom$\Gamma_1 \fCenter M : A \lor B$
\Axiom$x:A, \Gamma_2 \fCenter N_1 : C$
\Axiom$y:B, \Gamma_3 \fCenter N_2: C$
\RightLabel{\Elim{\lor}}
\TrinaryInf$\Gamma_1, \Gamma_2, \Gamma_3 \fCenter
\dcase{M}{x}{N_1}{y}{N_2} : C$
\DisplayProof
\end{tabular}
}
\\[4ex]
\multicolumn{2}{c}{\Axiom$\Gamma \fCenter N: \lfalse$
\RightLabel{\FalseInt}
\UnaryInf$\Gamma \fCenter \abort{A}{N} : A$
\DisplayProof}
\end{tabular}
\caption{विधानीय पद-चिन्हांकित अंतःप्रज्ञावादी नैसर्गिक निगमन
पद्धतीचे नियम \Log{tN2ip}.}
\label{int:pty:ter:tab:tN2ip}
\end{table}
















\section{निष्पत्तीचे सिद्धता-पदांत रूपांतर}\label{int:pty:pt:sec}

विधानीय नैसर्गिक निगमनातील सिद्धता चे सिद्धता-पदांत रूपांतर,
म्हणजे~\Log{N2i} मधील सिद्धता चे~\Log{tN2i}
मधील सिद्धता मध्ये भाषांतर, अगदी सरळ आहे.
निष्पत्तीमधील प्रत्येक क्रमवर्तीच्या उजव्या
बाजूच्या सूत्राच्या डावीकडे सिद्धता-पद लिहायचे;
म्हणजे $\Gamma \Sequent M : A$ या रूपातील
क्रमवर्ती मिळतात. मग संदर्भ~$\Gamma$ मध्ये~$M$ हे
$A$ ची \emph{साक्ष देते} असे म्हणू. कोणते सिद्धता-पद
लिहायचे हे पूर्णपणे \Log{tN2i} च्या नियमांनी ठरते.

\begin{prop}
विधानीय \Log{N2i} मधील $\Gamma \Sequent A$ ची प्रत्येक
सिद्धता~$\delta$ ही~\Log{tN2i} मधील
$\Gamma \Sequent N : A$ च्या सिद्धता~$\delta'$ शी
अनुरूप असते. सिद्धता-पद~$N$ हे~$\delta$ वरून
एकमेवपणे ठरते.
\end{prop}

\begin{proof}
~$\delta$ च्या उंचीवर विगमन करू. $\delta$ हे
$x:A \Sequent A$ हे स्वयंसिद्धक असेल, तर
$\delta'$ हे~$x:A \Sequent x:A$ हे स्वयंसिद्धक आहे.
$\delta$ एखाद्या अनुमान-नियमावर संपत असेल, तर
विगमन-गृहीतक प्रत्येक आधारविधानासाठी सिद्धता-पद आणि
ते पद असलेल्या क्रमवर्तीची \Log{tN2i}-सिद्धता देते.
मग \Log{tN2i} चा अनुरूप नियम लावतो;
\Log{tN2i} चा अनुमान-नियमच अनुरूप सिद्धता-पद निश्चित करतो.
\end{proof}
येथे दिलेल्या पद-नियमांची व्याप्ती विधानीय आहे;
संख्यापकांचे स्वतंत्र पद-संकारक परिभाषित केलेले नाहीत.

\begin{ex}
  \Log{N1i} मधील $(A \land B) \lif (B \land A)$
  ची पुढील अगदी साधी सिद्धता घ्या:
  \[
  \AxiomC{$\Discharge{A\land B}{x}$}
  \RightLabel{\Elim\land[2]}
  \UnaryInfC{$B$}
  \AxiomC{$\Discharge{A\land B}{x}$}
  \RightLabel{\Elim\land[1]}
  \UnaryInfC{$A$}
  \RightLabel{\Intro\land}
  \BinaryInfC{$B \land A$}
  \DischargeRule{\Intro\lif}{x}
  \UnaryInfC{$(A \land B) \lif (B \land A)$}
  \DisplayProof
  \]
  \Log{N2i} मध्ये ती अशी दिसते:
  \[
  \Axiom$x : A\land B \fCenter A\land B$
  \RightLabel{\Elim\land[2]}
  \UnaryInf$x : A\land B \fCenter B$
  \Axiom$x : A\land B \fCenter A\land B$
  \RightLabel{\Elim\land[1]}
  \UnaryInf$x : A\land B \fCenter A$
  \RightLabel{\Intro\land}
  \BinaryInf$x : A\land B \fCenter B \land A$
  \RightLabel{\Intro\lif}
  \UnaryInf$\fCenter (A \land B) \lif (B \land A)$
  \DisplayProof
  \]
  सिद्धता-पदे वरपासून खाली देऊ. स्वयंसिद्धकांची
  सिद्धता-पदे ही संदर्भातील~$x$ या चलासारखीच आहेत.
  मग $\Elim\land[i]$ शी संबंधित सिद्धता-पदे अनुक्रमे
  $\proj{2}{x}$ आणि~$\proj{1}{x}$ अशी ठरतात.
  \Intro{\land} लावताना ही दोन सिद्धता-पदे एकत्र करतो:
  $\pair{\proj{2}{x}}{\proj{1}{x}}$.
  शेवटी \Intro{\lif} नियमाने $\lambd$-अमूर्तीकरण
  केले जाते; त्यामुळे~$x$ बद्ध होतो आणि~$x$ ही खूण
  $A \land B$ या गृहीतकाला होती हे नोंदले जाते:
  $\lambd[\typeof{x}{A \land
  B}][\pair{\proj{2}{x}}{\proj{1}{x}}]$.
  मग \Log{tN2i} मधील सिद्धता अशी आहे:
  \[
  \Axiom$x : A\land B \fCenter x: A\land B$
  \RightLabel{\Elim\land[2]}
  \UnaryInf$x : A\land B \fCenter \proj{2}{x} : B$
  \Axiom$x : A\land B \fCenter x: A\land B$
  \RightLabel{\Elim\land[1]}
  \UnaryInf$x : A\land B \fCenter \proj{1}{x} : A$
  \RightLabel{\Intro\land}
  \BinaryInf$x : A\land B \fCenter \pair{\proj{2}{x}}{\proj{1}{x}} :
    B \land A$
  \RightLabel{\Intro\lif}
  \UnaryInf$\fCenter \lambd[\typeof{x}{A \land
  B}][\pair{\proj{2}{x}}{\proj{1}{x}}] : (A \land B) \lif (B \land A)$
  \DisplayProof
  \]
\end{ex}












\section{सिद्धता-पदांवरून निष्पत्ती पुन्हा मिळवणे}\label{int:pty:tp:sec}

आता उलट दिशा पाहू: सिद्धता-पद~$N$ वरून
\Log{N2i} मधील सिद्धता पुन्हा मिळवणे.
सर्वसाधारणपणे हे दिलेल्या सिद्धता-पदात मुक्त असलेल्या
प्रत्येक चल~$x$ साठी $x : A$ अशी जोडी धरणाऱ्या
संदर्भाच्या सापेक्षच करता येते. पण आधी मुक्त चले
नसलेले उदाहरण घेऊ, म्हणजे हे पद:
\[
\lambd[\typeof{x}{A \land
  B}][\pair{\proj{2}{x}}{\proj{1}{x}}]
\]
याचे रूप $\lambd[\typeof{x}{A \land B}][N_1]$ आहे.
अशा रूपाचे पद निर्माण करणारा \Log{tN2i} चा
एकमेव नियम~\Intro{\lif} असल्यामुळे~$N$ ने
सांकेतिक केलेली सिद्धता~\Intro{\lif} वर संपली
पाहिजे, म्हणजे तिचा शेवट असा आहे:
\[
  \Axiom$x : A\land B \fCenter \pair{\proj{2}{x}}{\proj{1}{x}} :
    C$
  \RightLabel{\Intro\lif}
  \UnaryInf$\fCenter \lambd[\typeof{x}{A \land
  B}][\pair{\proj{2}{x}}{\proj{1}{x}}] : (A \land B) \lif C$
  \DisplayProof
  \]
$C$ काय आहे हे अद्याप माहीत नाही; पण पूर्वांग
$A \land B$ आहे, आणि आधारविधानाच्या संदर्भात
$x:A \land B$ असलेच पाहिजे हे माहीत आहे.

आता आधारविधानाशी संबंधित सिद्धता-पद निश्चित झाले आहे:
$\pair{\proj{2}{x}}{\proj{1}{x}}$.
\emph{त्याने} सांकेतिक केलेली सिद्धता~\Intro{\land}
वर संपली पाहिजे. ~$C$ काय आहे हे अजून माहीत
नाही; पण तो~\Intro{\land} चा निष्कर्ष असल्यामुळे
संयोग असलाच पाहिजे, म्हणजे
$C \ident (C_1 \land C_2)$. पुनर्रचना पुढे चालू ठेवू:
\[
  \Axiom$x : A\land B \fCenter \proj{2}{x} : C_1$
  \Axiom$x : A\land B \fCenter \proj{1}{x} : C_2$
  \RightLabel{\Intro\land}
  \BinaryInf$x : A\land B \fCenter \pair{\proj{2}{x}}{\proj{1}{x}} :
    C_1 \land C_2$
  \RightLabel{\Intro\lif}
  \UnaryInf$\fCenter \lambd[\typeof{x}{A \land
  B}][\pair{\proj{2}{x}}{\proj{1}{x}}] : (A \land B) \lif (C_1 \land C_2)$
  \DisplayProof
  \]
$\pi_2(x)$ हे सिद्धता-पद सांगते की डावीकडील वरची
क्रमवर्ती~\Elim{\land}[2] ने मिळाली; म्हणजे तिचे
आधारविधान $D \land C_1$ या रूपाचे आहे. उजवीकडे
$C_2$ कडे नेणारे अनुमान~\Elim{\land}[1] असले
पाहिजे आणि त्याचे आधारविधान $C_2 \land E$ या
रूपाचे असले पाहिजे, हे ठरवता येते.
\[
  \Axiom$x : A\land B \fCenter x: D\land C_1$
  \RightLabel{\Elim\land[2]}
  \UnaryInf$x : A\land B \fCenter \proj{2}{x} : C_1$
  \Axiom$x : A\land B \fCenter x: C_2\land E$
  \RightLabel{\Elim\land[1]}
  \UnaryInf$x : A\land B \fCenter \proj{1}{x} : C_2$
  \RightLabel{\Intro\land}
  \BinaryInf$x : A\land B \fCenter \pair{\proj{2}{x}}{\proj{1}{x}} :
    C_1 \land C_2$
  \RightLabel{\Intro\lif}
  \UnaryInf$\fCenter \lambd[\typeof{x}{A \land
  B}][\pair{\proj{2}{x}}{\proj{1}{x}}] : (A \land B) \lif (C_1 \land C_2)$
  \DisplayProof
  \]
सर्वांत वरच्या क्रमवर्तींची सिद्धता-पदे आता चले
असल्यामुळे त्या स्वयंसिद्धकेच असल्या पाहिजेत.
त्यावरून $C_1$, $C_2$, $D$, आणि $E$ ठरतात:
$D \ident A$ आणि $C_1 \ident B$ नसतील,
तर डावी क्रमवर्ती स्वयंसिद्धक नाही; तसेच
$C_2 \ident A$ आणि $E \ident B$ नसतील,
तर उजवी क्रमवर्तीही स्वयंसिद्धक नाही.
अशा रीतीने संपूर्ण सिद्धता पुन्हा मिळाली:
\[
  \Axiom$x : A\land B \fCenter x: A\land B$
  \RightLabel{\Elim\land[2]}
  \UnaryInf$x : A\land B \fCenter \proj{2}{x} : B$
  \Axiom$x : A\land B \fCenter x: A\land B$
  \RightLabel{\Elim\land[1]}
  \UnaryInf$x : A\land B \fCenter \proj{1}{x} : A$
  \RightLabel{\Intro\land}
  \BinaryInf$x : A\land B \fCenter \pair{\proj{2}{x}}{\proj{1}{x}} :
    B \land A$
  \RightLabel{\Intro\lif}
  \UnaryInf$\fCenter \lambd[\typeof{x}{A \land
  B}][\pair{\proj{2}{x}}{\proj{1}{x}}] : (A \land B) \lif (B \land A)$
  \DisplayProof
  \]













\section{प्रकार}\label{int:pty:typ:sec}

$(MN)$ सारखे सिद्धता-पद स्वतःहून कोणत्या सिद्धतेचे
सिद्धता-पद आहे हे एकमेवपणे ठरवत नाही. कारण मुक्त चले,
म्हणजे स्वयंसिद्धकांची सिद्धता-पदे, संबंधित स्वयंसिद्धकांतील
सूत्रे सांगत नाहीत. पण ती सुसंगत संदर्भ~$\Gamma$ मध्ये
दिली आणि पद योग्य प्रकाराचे असेल, तर सिद्धता \emph{एकमेवपणे} ठरते.

\begin{defn}
सिद्धता-पद~$N$ साठी \emph{संदर्भ} म्हणजे प्रकार-खुणांचा
सुसंगत संच~$\Gamma$: कोणत्याही चलाला दोन वेगळे प्रकार
देता येत नाहीत, आणि~$N$ च्या प्रत्येक मुक्त चलासाठी
$\Gamma$ मध्ये $x:A$ अशी नेमकी एक जोडी असते.
\end{defn}

या व्याख्येनुसार~$\Gamma$ मध्ये $x:A$ ही जोडी
\emph{फक्त}~$N$ च्या मुक्त चलांसाठीच असावी अशी अट नाही.
उदाहरणार्थ, $\Gamma = \{x:A, y:B, z:C\}$ हा
$\pair{x}{y}$ साठी संदर्भ आहे.

\begin{defn}\label{int:pty:typ:defn:type}
  ~$\Gamma$ हा~$N$ साठी संदर्भ आहे असे माना.
  ~$\Gamma$ च्या सापेक्ष~$N$ चा \emph{प्रकार}
  पुढीलप्रमाणे विगमनाने परिभाषित करतो:
  \begin{enumerate}
  \item $x:A \in \Gamma$ असेल तेव्हाच~$x$ चा प्रकार~$A$ असतो.
  \item $x:A, \Gamma$ च्या सापेक्ष~$N$ चा प्रकार~$B$
    असेल, तर $\lambd[\typeof{x}{A}][N]$ चा प्रकार
    $A \lif B$ असतो.
  \item $N$ चा प्रकार $A \lif B$ आणि~$M$ चा प्रकार~$A$
    असेल, तर $(NM)$ चा प्रकार~$B$ असतो.
  \item $N$ चा प्रकार~$A$ आणि~$M$ चा प्रकार~$B$
    असेल, तर $\pair{N}{M}$ चा प्रकार $A \land B$ असतो.
  \item $N$ चा प्रकार~$A_1 \land A_2$ असेल, तर
    $\proj{i}{N}$ चा प्रकार~$A_i$ असतो.
  \item $N$ चा प्रकार~$A$ असेल, तर
    $\inj[B]{i}{N}$ चा प्रकार $i=1$ असताना
    $A \lor B$ आणि $i=2$ असताना $B \lor A$ असतो.
  \item $M$ चा प्रकार $A \lor B$, तसेच
    $x_1:A, \Gamma$ च्या सापेक्ष~$N_1$ चा प्रकार~$C$
    आणि $x_2:B, \Gamma$ च्या सापेक्ष~$N_2$ चा
    प्रकार~$C$ असेल, तर
    $\dcase{M}{x_1}{N_1}{x_2}{N_2}$ चा प्रकार~$C$ असतो.
  \item $N$ चा प्रकार~$\lfalse$ असेल, तर
    $\abort{A}{N}$ चा प्रकार~$A$ असतो.
  \end{enumerate}
\end{defn}
या नियमांत न बसणाऱ्या पदाला प्रकार नाही. अमूर्तीकरणाच्या
आणि प्रकरण-विभाजनाच्या बद्ध चलांना संदर्भातील चलांपासून वेगळी
नवी नावे देऊन वरील संदर्भ-वाढ करायची.

\begin{prop}
सिद्धता-पद~$N$ साठीच्या सुसंगत संदर्भ~$\Gamma$ च्या
सापेक्ष~$N$ ला जास्तीत जास्त एक प्रकार असतो. प्रकार
अस्तित्वात असेल, तर तो एकमेव असतो; प्रत्येक कच्च्या
सिद्धता-पदाला प्रकार मिळतोच असे नाही.
\end{prop}

\begin{proof}
  ~$N$ वर विगमनाने.
\end{proof}

संदर्भ~$\Gamma$ च्या सापेक्ष~$N$ चा प्रकार~$A$
असेल, तर $\Gamma \Sequent N : A$ असे लिहितो.
\ref{int:pty:typ:defn:type} मधील उपवाक्ये परिचित वाटली पाहिजेत:
त्यात एक लहानसा फरक वगळता \Log{tN2i} चेच नियम
आहेत—संदर्भ~$\Gamma$ सर्वत्र तोच राहतो. ही
व्याख्या \Log{tN3i} या कलनाच्या रूपात मांडता येते;
त्याचे नियम \ref{int:pty:typ:tab:tN3ip} मध्ये दिले आहेत.









\begin{table}
\begin{tabular}{@{}c@{}c@{}}
\multicolumn{2}{l}{\textbf{स्वयंसिद्धके:}}\\
\multicolumn{2}{c}{$\Gamma \Sequent x:A$ जेव्हा $x:A \in \Gamma$}\\[3ex]
\multicolumn{2}{l}{\textbf{अनुमान-नियम:}}\\[3ex]
\multicolumn{2}{c}{\Axiom$x: A, \Gamma \fCenter N:B$
\RightLabel{\Intro{\lif}}
\UnaryInf$\Gamma \fCenter \lambd[\typeof{x}{A}][N] : A \lif B $
\DisplayProof}
\\[4ex]
\multicolumn{2}{c}{\Axiom$\Gamma \fCenter N:A \lif B$
\Axiom$\Gamma \fCenter M:A$
\RightLabel{\Elim{\lif}}
\BinaryInf$\Gamma \fCenter (NM) : B$
\DisplayProof}
\\[4ex]
\multicolumn{2}{c}{\Axiom$\Gamma \fCenter N_1 : A$
\Axiom$\Gamma \fCenter N_2 : B$
\RightLabel{\Intro{\land}}
\BinaryInf$\Gamma \fCenter \pair{N_1}{N_2} : A \land B $
\DisplayProof}
\\[4ex]
\Axiom$\Gamma \fCenter N: A \land B$
\RightLabel{\Elim{\land}[1]}
\UnaryInf$\Gamma \fCenter \proj{1}{N} : A$
\DisplayProof
&
\Axiom$\Gamma \fCenter N:A \land B$
\RightLabel{\Elim{\land}[2]}
\UnaryInf$\Gamma \fCenter \proj{2}{N} : B$
\DisplayProof\\[3ex]
\\[4ex]
\Axiom$\Gamma \fCenter N: A$
\RightLabel{\Intro{\lor}[1]}
\UnaryInf$\Gamma \fCenter \inj[B]{1}{N} : A \lor B$
\DisplayProof
&
\Axiom$ \Gamma \fCenter N : A$
\RightLabel{\Intro{\lor}[2]}
\UnaryInf$\Gamma \fCenter \inj[B]{2}{N} : B \lor A$
\DisplayProof\\[3ex]
\multicolumn{2}{c}{
\begin{tabular}[b]{@{}l@{}}
\Axiom$\Gamma \fCenter M : A \lor B$
\Axiom$x:A, \Gamma \fCenter N_1 : C$
\Axiom$y:B, \Gamma \fCenter N_2: C$
\RightLabel{\Elim{\lor}}
\TrinaryInf$\Gamma \fCenter
\dcase{M}{x}{N_1}{y}{N_2} : C$
\DisplayProof
\end{tabular}
}
\\[4ex]
\multicolumn{2}{c}{\Axiom$\Gamma \fCenter N: \lfalse$
\RightLabel{\FalseInt}
\UnaryInf$\Gamma \fCenter \abort{A}{N} : A$
\DisplayProof}
\end{tabular}
\caption{प्रकार-निर्देशन पद्धती म्हणून वापरल्या जाणाऱ्या
\Log{tN3ip} या विधानीय पद-चिन्हांकित अंतःप्रज्ञावादी
नैसर्गिक निगमन पद्धतीचे नियम.}
\label{int:pty:typ:tab:tN3ip}
\end{table}




सिद्धता-पदे ही $\lambd$-कलनाच्या एका आवृत्तीतील
पदे आहेत: \emph{गुणाकार, बेरीज आणि रिक्त प्रकार
असलेले प्रकारयुक्त $\lambd$-कलन}. पारंपरिक
$\lambd$-कलनात फक्त चलांपासून बनलेली पदे,
\emph{$\lambd$-अमूर्तीकरणे}~$\lambd[x][N]$, आणि
उपयोजन $(MN)$ असते; त्यात प्रकारयुक्त
$\lambd$-अमूर्तीकरण~$\lambd[\typeof{x}{A}][N]$
यातील~$A$ हे प्रकार-चिन्ह नसते. ~$\lambd$-कलन
हे ट्यूरिंग-संपूर्ण संगणन-प्रतिमान आहे (म्हणजे प्रत्येक
आंशिक संगणनीय फलन त्यातील पदाने दर्शवता येते).
सिद्धता-पदे आणि प्रकार-नियम देतात ती~$\lambd$-कलनाची
आवृत्ती निर्बंधित संगणन-प्रतिमान आहे; पण तिच्या मूलभूत
कल्पना त्याच आहेत. हे निर्बंध प्रकार-निर्देशनातून येतात:
फक्त \emph{योग्य प्रकाराची} पदे, म्हणजे
संदर्भ~$\Gamma$ च्या सापेक्ष प्रकार असलेली
सिद्धता-पदे, ग्राह्य असतात. उदाहरणार्थ,
$\lambd[x][(xx)]$ हे पारंपरिक प्रकारविरहित
$\lambd$-कलनात ग्राह्य पद आहे; पण
$\lambd[\typeof{x}{A}][(xx)]$ हे कोणत्याही~$A$
साठी योग्य प्रकाराचे नाही. कारण $(xx)$ ला प्रकार
असण्यासाठी पहिल्या~$x$ चा प्रकार $A \lif B$
आणि दुसऱ्या~$x$ चा प्रकार~$A$ असावा लागेल.
पण~$A$ हे $A \lif B$ चे योग्य उप-सूत्र
असल्यामुळे हे अशक्य आहे.

एखाद्या प्रकाराचे पद ज्या वस्तूला नाव देते, तिचा
प्रकार म्हणून प्रकारांकडे सहजपणे पाहता येते. म्हणून
$A \land B$ या प्रकाराचे पद पहिला घटक~$A$
आणि दुसरा~$B$ या प्रकाराचा असलेली जोडी निवडते;
म्हणजे गुणाकार~$A \times B$ मधील घटक.
$A \lif B$ या प्रकाराचे पद प्रकार~$A$ च्या
वस्तूंपासून प्रकार~$B$ च्या वस्तूंकडे जाणारे फलन आहे.
$A \lor B$ या प्रकारचे पद हे प्रकार~$A$ किंवा
प्रकार~$B$ ची वस्तू असते; तसेच ती दोहोंपैकी कोणत्या
प्रकाराची आहे ही माहिती त्यात असते. हे~$A$ आणि~$B$
या दोन संचांच्या असंयुक्त संघासारखे आहे:
$\{\tuple{i,x} : i = 1 \text{ किंवा } 2, x \in A
\text{ जर } i = 1, x \in B \text{ जर } i=2\}$.
$\lfalse$ हा प्रकार रिक्त संच आहे; म्हणजे त्या प्रकाराचे
पद कोणत्याही वस्तूला नाव देऊ शकत नाही.
नैसर्गिक निगमन (\Log{tN3i} सह) सुसंगत असल्यामुळे
खुली गृहीतके नसताना~$\lfalse$ ची सिद्धता नसते.
म्हणून~$\lfalse$ प्रकाराचे बंद सिद्धता-पद
(मुक्त चले नसलेले पद) नसते.

प्रकारयुक्त $\lambd$-कलनाचे संकारक योग्य प्रकाराच्या
वस्तूंना सहजपणे नाव देणारी पदे तयार करतात, अर्थात
त्यांना लावलेली पदेही ठरावीक प्रकारच्या वस्तूंना नाव
देत असतील तर. उदाहरणार्थ, ``$N_1 : A$ आणि
$N_2 : B$ असतील, तर
$\pair{N_1}{N_2} : A \land B$'' याचा अर्थ:
$N_1$ ने प्रकार~$A$ च्या वस्तूला आणि~$N_2$ ने
प्रकार~$B$ च्या वस्तूला नाव दिले, तर
$\pair{N_1}{N_2}$ हे प्रकार $A \land B$ च्या
वस्तूला नाव देते.













\section{न्यूनीकरण}\label{int:pty:red:sec}

नैसर्गिक निगमनाच्या निष्पत्तीमध्ये
परिचय नियमापाठोपाठ निरसन नियम
आला, तर तो वळसा ``न्यूनीत'' करता येतो. उदाहरणार्थ,
$\Intro{\land}$ नंतर $\Elim{\land}$ आले, तर ते
``एकमेकांना रद्द'' करतात: $\Elim{\land}$ चा निष्कर्ष
नेहमीच \Intro{\land} ने तयार केलेल्या संयोगातील
एका घटकासारखा असतो, आणि $\Intro{\land}$ ची दोन्ही
आधारविधाने तेच संयोगघटक असतात. म्हणून एखाद्या
संयोगघटकाची निष्पत्ती हवी असेल, तर संयोग
प्रस्तुत करून त्याचे निरसन न करता तीच निष्पत्ती
थेट वापरता येते. आपल्याला असे मिळते:
\begin{gather*}
  \bottomAlignProof
  \AxiomC{}
  \RightLabel{$\delta_1$}
  \DeduceC{$A_1$}
  \AxiomC{}
  \RightLabel{$\delta_2$}
  \DeduceC{$A_2$}
  \RightLabel{$\Intro{\land}$}
  \BinaryInfC{$A_1 \land A_2$}
  \RightLabel{$\Elim{\land}_i$}
  \UnaryInfC{$A_i$}
  \DisplayProof\bottomAlignProof
  \quad
  \redone
  \quad
  \AxiomC{}
  \RightLabel{$\delta_i$}
  \DeduceC{$A_i$}
  \DisplayProof
\end{gather*}

या सोपीकरणातून संबंधित सिद्धता-पदांवरील तसाच
न्यूनीकरण-संबंध दिसतो. ~$N_1$ हे~$\delta_1$ चे
आणि~$N_2$ हे~$\delta_2$ चे सिद्धता-पद असेल,
तर डावीकडील सिद्धतेचे सिद्धता-पद एका पायरीत
उजवीकडील सिद्धतेच्या सिद्धता-पदात न्यूनीत होते:
\[
  \proj{i}{\pair{N_1}{N_2}} \redone N_i
\]
प्रकारयुक्त लॅम्डा कलनात हा गुणाकार-प्रकाराचा
\emph{बीटा-न्यूनीकरण} नियम आहे.

नैसर्गिक निगमनात \Intro{\land} नंतर
\Elim{\land} अशी अनुमानांची जोडी, म्हणजे न्यूनीकरण
लागू होणारी जोडी, हिला \emph{कट} म्हणतात. प्रकारयुक्त
लॅम्डा कलनात $\redone$ च्या डावीकडील पदाला
\emph{रेडेक्स}, तर उजवीकडील पदाला \emph{संकुचनफल}
म्हणतात. प्रकारविरहित लॅम्डा कलनात फक्त
$(\lambd[x][N])Q$ यालाच रेडेक्स मानतात; पण
प्रकारयुक्त लॅम्डा कलनात $\pair{N}{M}$,
$\proj{i}{N}$, $\inj[A]{i}{N}$,
$\dcase{N}{x_1}{M_1}{x_2}{M_2}$, आणि
$\abort{A}{N}$ यांसारखी पदेही विन्यासात येतात,
आणि त्यांना अनुरूप रेडेक्सही असतात.

त्याचप्रमाणे वियोगासाठीही न्यूनीकरण आहे:
\begin{gather*}
  \bottomAlignProof
  \AxiomC{}
  \RightLabel{$\delta$}
  \DeduceC{$A_i$}
  \RightLabel{$\Intro{\lor}$}
  \UnaryInfC{$A_1 \lor A_2$}
  \AxiomC{$\Discharge{A_1}{x_1}$}
  \RightLabel{$\delta_1$}
  \DeduceC{$C$}
  \AxiomC{$\Discharge{A_2}{x_2}$}
  \RightLabel{$\delta_2$}
  \DeduceC{$C$}
  \DischargeRule{$\Elim{\lor}$}{x_1,x_2}
  \TrinaryInfC{$C$}
  \DisplayProof\bottomAlignProof
  \quad
  \redone
  \quad
  \AxiomC{}
  \RightLabel{$\delta$}
  \DeduceC{$A_i$}
  \RightLabel{$\delta_i$}
  \DeduceC{$C$}
  \DisplayProof
\end{gather*}
याला अनुरूप सिद्धता-पदांवरील न्यूनीकरण असे आहे:
\begin{gather*}
  \dcase{\inj[A_{3-i}]{i}{M}}{x_1}{N_1}{x_2}{N_2} \redone \Subst{N_i}{M}{x_i}
\end{gather*}
येथे~$M$ हे~$\delta$ चे, आणि~$N_i$ हे~$\delta_i$
चे सिद्धता-पद आहे. हा बेरीज-प्रकारांसाठीचा
बीटा-न्यूनीकरण नियम आहे. येथे
$\Subst{M}{N_i}{x}$ म्हणजे~$M$ मधील~$x$ या
चराची सर्व मुक्त स्थाने~$N_i$ ने बदलणे.
हे मुक्त-चल प्रतिस्थापन बद्ध-चल ग्रहण होऊ न देता
करायचे; आवश्यक बद्ध चलांना आधी नवी नावे द्यायची.
विन्यासातील प्रत्येक संकारकालाच रेडेक्स म्हणत नाही:
मुख्य बीटा-रेडेक्स हे वरच्या गुणाकाराच्या आणि पुढच्या
बेरीज व फलनाच्या विशिष्ट रूपांचे असतात.

फलन-प्रकाराचे न्यूनीकरण म्हणजे $\Intro\lif$ नंतर
$\Elim\lif$ आल्याने झालेला वळसा काढणे:
\begin{gather*}
  \bottomAlignProof
  \AxiomC{$\Discharge{A}{x}$}
  \RightLabel{$\delta$}
  \DeduceC{$B$}
  \DischargeRule{$\Intro{\lif}$}{x}
  \UnaryInfC{$A \lif B$}
  \AxiomC{}
  \RightLabel{$\delta'$}
  \DeduceC{$A$}
  \RightLabel{$\Elim{\lif}$}
  \BinaryInfC{$B$}
  \DisplayProof\bottomAlignProof
  \quad
  \redone
  \quad
  \AxiomC{}
  \RightLabel{$\delta'$}
  \DeduceC{$A$}
  \RightLabel{$\delta$}
  \DeduceC{$B$}
  \DisplayProof
\end{gather*}
सिद्धता-पदांसाठी हे नेहमीचे बीटा-न्यूनीकरण आहे:
\begin{gather*}
  (\lambd[\typeof{x}{A}][N]M)
  \redone \Subst{N}{M}{x}
\end{gather*}

सिद्धता वरील न्यूनीकरण रूपांतरणांखेरीज
क्रमपरिवर्तन रूपांतरणांचाही विचार करता येतो.
ती \Elim{\lor} नंतर दुसरा निरसन नियम
असलेल्या (उप-)सिद्धता वर लागू होतात, उदाहरणार्थ:
\begin{multline*}
  \bottomAlignProof
  \AxiomC{}
  \RightLabel{$\delta$}
  \DeduceC{$A_i$}
  \RightLabel{$\Intro{\lor}$}
  \UnaryInfC{$A_1 \lor A_2$}
  \AxiomC{$\Discharge{A_1}{x_1}$}
  \RightLabel{$\delta_1$}
  \DeduceC{$C_1 \land C_2$}
  \AxiomC{$\Discharge{A_2}{x_2}$}
  \RightLabel{$\delta_2$}
  \DeduceC{$C_1 \land C_2$}
  \DischargeRule{$\Elim{\lor}$}{x_1,x_2}
  \TrinaryInfC{$C_1 \land C_2$}
  \RightLabel{\Elim{\land}[i]}
  \UnaryInfC{$C_i$}
  \DisplayProof\bottomAlignProof
  \quad
  \redone
  \\
  \AxiomC{}
  \RightLabel{$\delta$}
  \DeduceC{$A_i$}
  \RightLabel{$\Intro{\lor}$}
  \UnaryInfC{$A_1 \lor A_2$}
  \AxiomC{$\Discharge{A_1}{x_1}$}
  \RightLabel{$\delta_1$}
  \DeduceC{$C_1 \land C_2$}
  \RightLabel{\Elim{\land}[i]}
  \UnaryInfC{$C_i$}
  \AxiomC{$\Discharge{A_2}{x_2}$}
  \RightLabel{$\delta_2$}
  \DeduceC{$C_1 \land C_2$}
  \RightLabel{\Elim{\land}[i]}
  \UnaryInfC{$C_i$}
  \DischargeRule{$\Elim{\lor}$}{x_1,x_2}
  \TrinaryInfC{$C_i$}
  \DisplayProof
\end{multline*}
~$\delta$, $\delta_1$, आणि~$\delta_2$ यांची
सिद्धता-पदे अनुक्रमे~$M$, $N_1$, आणि~$N_2$
असतील, तर सिद्धता-पदांसाठी, म्हणजेच~$\lambd$
पदांसाठी, अनुरूप न्यूनीकरण नियम असा आहे:
\[
  \proj{i}{\dcase{M}{x_1}{N_1}{x_2}{N_2}} \redone
\dcase{M}{x_1}{\proj{i}{N_1}}{x_2}{\proj{i}{N_2}}
\]
इतर निरसनांसाठीही वियोग-निरसनाच्या निष्कर्षावरील
निरसन दोन्ही शाखांच्या निष्कर्षांवर हलवायचे. बाहेरील
आदान किंवा उपपद आत जाताना~$x_1,x_2$ यांनी त्याचे
मुक्त चल पकडू नयेत म्हणून या बद्ध चलांना आधी नवी
नावे द्यायची. असत्याच्या सरलीकरणांसाठी नैसर्गिक
निगमनाच्या सामान्यरूपीकरण प्रकरणातील अनुरूप
रूपांतरणे वापरायची.

अर्थात, वळसे केवळ सिद्धता च्या शेवटीच नव्हे,
तर सिद्धता च्या आतही काढता येतात: वळशावर संपणाऱ्या
उप-सिद्धता ला न्यूनीकरण रूपांतरण लावायचे.
तसेच $\lambd$-पदांच्या उपपदांवर $\beta$-न्यूनीकरण
लागू होते. $N$ हे~$M$ चे उपपद असेल आणि
$N \redone N'$ असेल, तर~$M$ मधील उपपद~$N$ हे~$N'$ ने
बदलून~$M'$ मिळवता येते. अशा वेळीही $M \redone M'$
असे लिहितो. $M = M_1 \redone M_2 \redone M_3
\redone \dots \redone M_k = M'$ असा क्रम असेल
(यात $k=1$, म्हणजे $M = M'$ ही बाबही येते),
तर $M \red M'$ असे लिहितो.

ज्या सिद्धता च्या कोणत्याही उप-सिद्धता वर
रूपांतरण लागू होत नाही, तिला \emph{सामान्य} म्हणतात.
तसेच ज्या $\lambd$-पदावर (त्याच्या कोणत्याही उपपदावर)
रूपांतरण लागू होत नाही, त्याला \emph{सामान्य} म्हणतात.
सामान्य $\lambd$-पदाला \emph{सामान्य रूप} असेही म्हणतात.














\section{प्रकार-संरक्षण}\label{int:pty:tpr:sec}

विचारात घेतलेली न्यूनीकरण आणि क्रमपरिवर्तन रूपांतरणे
\Log{tN2i} आणि \Log{tN3i} या पद्धतींसाठीही थेट
परिभाषित करता येतात. \Log{tN3i} पद्धत संदर्भ~$\Gamma$
च्या सापेक्ष $\lambd$-पद~$N$ चा प्रकार असे
सूत्र~$A$ म्हणून परिभाषित करते की
$\Gamma \Sequent N : A$ ची \Log{tN3i} मध्ये
सिद्धता आहे. $\lambd$-पदांसाठीची
पद-रूपांतरणे सिद्धता ची न्यूनीकरण
आणि क्रमपरिवर्तन रूपांतरणे तंतोतंत प्रतिबिंबित करतात.
त्याचा महत्त्वाचा परिणाम म्हणजे संदर्भ~$\Gamma$ च्या
सापेक्ष $\lambd$-पदाचा प्रकार अनुरूप रूपांतरणानंतरही
तोच राहतो. या गुणधर्माला \emph{प्रकार-संरक्षण} म्हणतात.

\begin{prop}
संदर्भ~$\Gamma$ च्या सापेक्ष $N$ चा प्रकार~$A$
असेल आणि $N \redone N'$ असेल, तर~$\Gamma$ च्याच
सापेक्ष $N'$ चा प्रकारही~$A$ असतो.
\end{prop}

\begin{proof}
  ~$N$ वर आणि $\Gamma \Sequent N : A$ च्या
  सिद्धता~$\delta$ च्या उंचीवर विगमन करून पुढील
  गोष्ट सहज सिद्ध करता येते:~$M$ हे~$N$ चे उपपद
  असेल, तर~$\delta$ मध्ये $\Gamma' \Sequent M : B$
  वर संपणारी उप-सिद्धता~$\delta_1$ असते.
  $M$ हा रेडेक्स असेल, तर~$\delta_1$ वर रूपांतरण
  लागू होते. ~$\delta_1$ ला न्यूनीकरण लावल्यावर
  $\Gamma' \Sequent M' : B$ ची सिद्धता~$\delta_1'$
  मिळते. \Log{tN3i} च्या रूपांतरणांची तपासणी केली,
  तर~$M'$ हे~$M$ चे संकुचनफल आहे. ~$\delta$ मध्ये
  $\delta_1$ च्या जागी~$\delta_1'$ ठेवल्यावर
  $\Gamma \Sequent N' : A$ ची सिद्धता~$\delta'$
  मिळते; येथे~$N$ मधील उपपद~$M$ च्या जागी~$M'$
  ठेवल्यावर~$N'$ मिळते.

  उदाहरणार्थ, \Intro{\land} नंतर
  \Elim{\land} यासाठीचे न्यूनीकरण रूपांतरण घ्या:
  \[
  \bottomAlignProof
  \AxiomC{}
  \Deduce$\Gamma \fCenter N_1 : B_1$
  \AxiomC{}
  \Deduce$\Gamma \fCenter N_2 : B_2$
  \RightLabel{\Intro{\land}}
  \BinaryInf$\Gamma \fCenter \pair{N_1}{N_2} : B_1 \land B_2$
  \RightLabel{\Elim{\land}[i]}
  \UnaryInf$\Gamma \fCenter \proj{i}{\pair{N_1}{N_2}} : B_i$
  \DisplayProof
  \quad\redone\quad
  \bottomAlignProof
  \AxiomC{}
  \Deduce$\Gamma \fCenter N_i : B_i$
  \DisplayProof
  \]
  उजवीकडील सिद्धता चे सिद्धता-पद, म्हणजे~$\delta_1'$
  हे~$N_i$ आहे. डावीकडील~($\delta_1$) सिद्धता
  च्या $\proj{i}{\pair{N_1}{N_2}}$ या सिद्धता-पदावर
  $\beta$-परिवर्तन लावल्यावर मिळणारे हेच फल आहे.

  किंवा \Intro{\lif} नंतर \Elim{\lif} येणारे
  अनुमान आणि त्यावर लागू होणारे न्यूनीकरण रूपांतरण घ्या:
  \begin{multline*}
  \bottomAlignProof
  \AxiomC{$\delta_2$}
  \Deduce$x: B_1, \Gamma \fCenter N : B_2$
  \RightLabel{\Intro{\lif}}
  \UnaryInf$\Gamma \fCenter \lambd[\typeof{x}{B_1}][N] : B_1 \lif B_2$
  \AxiomC{}
  \RightLabel{$\delta_3$}
  \Deduce$\Gamma \fCenter M : B_1$
  \RightLabel{\Elim{\lif}}
  \BinaryInf$\Gamma \fCenter (\lambd[\typeof{x}{B_1}][N] M) : B_2$
  \DisplayProof
  \quad\redone\\
  \bottomAlignProof
  \AxiomC{}
  \RightLabel{$\delta_3$}
  \Deduce$\Gamma \fCenter M : B_1$
  \RightLabel{$\Subst{\delta_2}{\delta_3}{x:B_1}$}
  \Deduce$\Gamma \fCenter \Subst{N}{M}{x} : B_2$
  \DisplayProof
  \end{multline*}
  पुन्हा उजवीकडील सिद्धता चे सिद्धता-पद हे
  डावीकडील सिद्धता च्या सिद्धता-पदाचे
  $\beta$-न्यूनीकरण केल्यावर मिळणारेच फल आहे.
  अर्थात, ~$\delta_2$ च्या उंचीवर विगमन करून पुढील
  गोष्ट काळजीपूर्वक तपासावी लागेल:~$\delta_2$
  मधील $\Gamma' \Sequent x:B_1$ या रूपातील सर्व
  स्वयंसिद्धकांच्या जागी $\Gamma \Sequent M:B_1$
  ची सिद्धता~$\delta_3$ ठेवल्यास खरोखर
  $\Gamma \Sequent \Subst{N}{M}{x} : B_2$
  ची सिद्धता~$\Subst{\delta_2}{\delta_3}{x:B_1}$
  मिळते.
\end{proof}
हे प्रतिस्थापन करताना बद्ध चलांना आधी सुरक्षित नवी
नावे द्यायची आणि संबंधित खुणांनी बांधलेली स्वयंसिद्धकेच
बदलायची. येथे संदर्भ वाढवता येणारी \Log{tN3ip}
पद्धत वापरतो. \Log{tN2ip} मध्ये वापरलेल्या गृहीतकांचा
अचूक संच नोंदतो; शाखा किंवा न वापरलेले आदान गळल्यास
तो संच कमी होऊ शकतो, म्हणून तिथे तोच अचूक संदर्भ
राहतो असा दावा नाही.

सिद्धता आणि $\lambd$-पदे यांतील निकटच्या
अनुरूपतेचा आणखी एक परिणाम आहे. संदर्भ~$\Gamma$
च्या सापेक्ष~$A$ या प्रकाराचे $\lambd$-पद~$N$
मध्ये रेडेक्स~$M$ असेल (म्हणजे ते सामान्य रूपात
नसेल), तर $\Gamma \Sequent N : A$ ची अनुरूप
\Log{tN3i} सिद्धता~$\delta$ यात
$\Gamma' \Sequent M: B$ वर संपणारी अशी
उप-सिद्धता असते की तिच्यावर रूपांतरण लागू होते.
संकुचनफलाने~$M$ बदलून $N \redone N'$ होत असेल,
तर~$\delta$ ला न्यूनीकरण लावून
$\Gamma \Sequent N' : A$ ची सिद्धता मिळते.
म्हणून सामान्य रूपात नसलेले प्रत्येक $\lambd$-पद
दुसऱ्या पदात अनुरूप रूपांतरणाने न्यूनीत होते. या गुणधर्माला
\emph{प्रगती} म्हणतात.
येथील ‘सामान्य’ याचा अर्थ घोषित रूपांतरण लागू
नसलेले पद असा आहे. खुले सामान्य पद हे मूल्याच्या
नेहमीच्या संगणकीय रूपात असतेच, किंवा प्रत्येक टप्प्यावर
फक्त मुख्य बीटा-पायरी लागू होतेच, असा दावा नाही.

प्रगती आणि प्रकार-संरक्षण मिळून $\lambd$-पदांचे
घोषित पद-न्यूनीकरण ``अडकत'' नाही याची खात्री देतात:
एखादे पद योग्य प्रकाराचे असेल आणि सामान्य रूपात
नसेल, तर ते दुसऱ्या योग्य प्रकाराच्या पदात न्यूनीत
होते (आणि त्या पदाचा प्रकारही तोच असतो).














\section{सामान्यीकरण}\label{int:pty:nor:sec}

या विभागात योग्य प्रकाराच्या सिद्धता-पदांना सामान्य रूप
मिळते हे दाखवू. यासाठी आधी सिद्ध केलेले नैसर्गिक
निगमनाचे सामान्यरूपीकरण प्रमेय आणि विधाने-हे-प्रकार
अनुरूपता वापरू. येथे पूर्ण सामान्यरूपीकरणासाठी
न्यूनीकरणासह आवश्यक क्रमपरिवर्तन आणि असत्याच्या
सरलीकरणाचे अनुरूप पद-रूपांतरणही धरतो; फक्त तीन
मुख्य $\beta$-नियम पुरेसे असल्याचा दावा नाही.

प्रथम पदांची गुंतागुंत मोजणारी काही फलने परिभाषित करू.
सूत्रांची \emph{लांबी}~$\len{A}$ पुढीलप्रमाणे:
\begin{align*}
  \len{p} &= 0\\
  \len{A \land B} &= \len{A} + \len{B} + 1\\
  \len{A \lor B} &= \len{A} + \len{B} + 1 \\
  \len{A \lif B} &= \len{A} + \len{B} + 1.
\end{align*}
मुख्य $\beta$-रेडेक्स~$M$ ची गुंतागुंत त्याच्या \emph{कट-कोटी}ने,
म्हणजे~$\cutrank{M}$ ने मोजतो:
\begin{align*}
  \cutrank{(\lambd[\typeof{x}{A}][\typeof{N}{B}])Q} &=
  \len{A} + \len{B} + 1 \\
  \cutrank{\ande{i}{\andi{\typeof{M}{A}}{\typeof{N}{B}}}} &=
  \len{A} + \len{B} + 1 \\
  \cutrank{\ore{\inj[A_{3-i}]{i}{\typeof{M}{A_i}}}
    {\typeof{x_1}{A_1}}{\typeof{N_1}{C}}
    {\typeof{x_2}{A_2}}{\typeof{N_2}{C}}} &=
  \len{A_1} + \len{A_2} + 1
\end{align*}
या तीन मुख्य बीटा-प्रकारांसाठी सिद्धता-पदाची
गुंतागुंत त्यातील सर्वाधिक गुंतागुंतीच्या रेडेक्सने मोजतो;
असे रेडेक्स नसतील, तर ती~$0$ असते. हा मापदंड
क्रमपरिवर्तन किंवा असत्याच्या सरलीकरणाचे मोजमाप नाही:
\[
  \maxrank{M} = \max \{\cutrank{N} | N \text{ हे उपपद आहे } M \text{ आणि रेडेक्स आहे}\}
\]
अणुरूप प्रकार आणि~$\lfalse$ यांची लांबी~$0$ धरतो.
वरील महत्तम मूल्य घेण्यासाठीचा संच रिकामा असेल, तर व्याख्येने
$\maxrank{M}=0$ घेतो.
बेरीज-प्रकाराच्या पदात~$\inj[A_{3-i}]{i}{M}$ वापरतो:
$M$ चा प्रकार~$A_i$ आहे आणि वरची प्रकार-खूण
निवडला न गेलेला दुसरा प्रकार सांगते. त्यामुळे पूर्ण
प्रकार~$A_1\lor A_2$ मिळतो, प्रकार-तक्त्याप्रमाणे.

\begin{lem}\label{int:pty:nor:lem:subst}
  $\Subst{M}{\typeof{N}{A}}{\typeof{x}{A}}$ हा मुख्य बीटा-रेडेक्स असेल
  आणि $M \not\ident x$ असेल, तर पुढीलपैकी एक शक्यता असते:
  \begin{enumerate}
  \item $M$ स्वतः रेडेक्स आहे; किंवा
  \item $M$ चे रूप $\ande{i}{x}$ आणि $N$ चे रूप
    $\andi{P_1}{P_2}$ आहे;
  \item $M$ चे रूप $\ore{x}{x_1}{P_1}{x_2}{P_2}$ आणि $N$ चे
    रूप $\inj{i}{Q}$ आहे;
  \item $M$ चे रूप $x Q$ आणि $N$ चे रूप
    $\lambd[x][P]$ आहे.
  \end{enumerate}
  पहिल्या बाबतीत $\cutrank{\Subst{M}{N}{x}} = \cutrank{M}$;
  इतर बाबतींत $\cutrank{\Subst{M}{N}{x}} = \len{A}$.
\end{lem}
\begin{proof}
  ~$M$ वर विगमनाने सिद्ध करू.
  \begin{enumerate}
  \item $M$ हे एकच चर $y$ असेल आणि $y \not\ident x$ असेल,
    तर $\Subst{y}{N}{x}$ हे $y$ आहे; म्हणून ते रेडेक्स नाही.
  \item $M$ चे रूप $\andi{N_1}{N_2}$, किंवा $\lambd[x][N]$,
    किंवा $\ori{i}{A}{N}$ असेल, तर
    $\Subst{M}{\typeof{N}{A}}{\typeof{x}{A}}$ चेही तेच
    रूप असते; म्हणून ते रेडेक्स नसते.
  \item $M$ चे रूप $\ande{i}{P}$ असेल, तर दोन बाबी पाहू.
    \begin{enumerate}
      \item $P$ चे रूप $\andi{P_1}{P_2}$ असेल, तर
        $M \ident \ande{i}{\andi{P_1}{P_2}}$ हा रेडेक्स आहे;
        आणि स्पष्टच,
        \[
        \Subst{M}{N}{x} \ident
        \ande{i}{\andi{\Subst{P_1}{N}{x}}{\Subst{P_2}{N}{x}}}
        \]
        हाही रेडेक्स आहे. दोघांच्या कट-कोटी समान आहेत.
      \item $P$ एकच चर असेल, तर प्रतिस्थापनाने रेडेक्स
        मिळण्यासाठी तो~$x$ असलाच पाहिजे, आणि~$N$ चे रूप
        $\andi{P_1}{P_2}$ असले पाहिजे. आता
        \[
        \Subst{M}{N}{x} \ident \Subst{\ande{i}{x}}{\andi{P_1}{P_2}}{x},
        \]
        विचारात घ्या. हे $\ande{i}{\andi{P_1}{P_2}}$ आहे.
        त्याची कट-कोटी $\cutrank{\ande{i}{\andi{P_1}{P_2}}}$
        म्हणजेच $\len{A}$ आहे.
    \end{enumerate}
  \end{enumerate}
  $\ore{N}{x_1}{N_1}{x_2}{N_2}$ आणि $P Q$ या बाबीही
  अशाच प्रकारच्या आहेत.
\end{proof}
येथील प्रतिस्थापन मुक्त चलांना बद्ध होऊ न देता करायचे;
आवश्यकतेनुसार बद्ध चलांना आधी नवी नावे द्यायची.
नवीन रेडेक्सच्या वर्गीकरणात बदललेल्या पदाच्या
मुळाशी असलेला रेडेक्स विचारात घेतला आहे, सर्व
आतील रेडेक्स नव्हेत. $\abort{B}{P}$ चे मूळ
तीन मुख्य बीटा-रेडेक्सपैकी कोणतेही नसते.

\begin{lem}
  $M$ या मुख्य बीटा-रेडेक्सचे $M'$ मध्ये संकुचन होत असेल आणि~$M$ च्या प्रत्येक
  योग्य रेडेक्स उपपद~$N$ साठी
  $\cutrank{M} > \cutrank{N}$ असेल, तर
  $\cutrank{M} > \maxrank{M'}$.
\end{lem}

\begin{proof}
  बाबींनुसार सिद्ध करू.
  \begin{enumerate}
  \item $M$ चे रूप $\ande{i}{\andi{M_1}{M_2}}$ असेल, तर
    $M'$ हे $M_i$ आहे. $M_i$ चे कोणतेही उपपद हे~$M$ चे
    योग्य उपपदही असल्यामुळे दावा सिद्ध होतो.
  \item $M$ चे रूप
    $(\lambd[\typeof{x}{A}][N])\typeof{Q}{A}$ असेल,
    तर $M'$ हे $\Subst{N}{\typeof{Q}{A}}{\typeof{x}{A}}$
    आहे. $M'$ मधील रेडेक्स घ्या. एकतर तो~$Q$ च्या बसवलेल्या प्रतीत आहे; अशा वेळी
    त्याची कट-कोटी गृहीतकामुळे मूळ~$M$ पेक्षा कमी आहे.
    अन्यथा~$N$ मध्ये त्याला
    अनुरूप व समान कट-कोटीचा रेडेक्स आहे; गृहीतकामुळे ती
    कट-कोटी $\cutrank{M}$ पेक्षा कमी आहे. अन्यथा त्याची
    कट-कोटी $\len{A}$ आहे; व्याख्येनुसार ती
    $\cutrank{(\lambd[x^{A}][N])Q}$ पेक्षा कमी आहे.
  \item $M$ चे रूप
    \[
    \ore{\inj[A_{3-i}]{i}{\typeof{N}{A_i}}}
        {\typeof{x_1}{A_1}}{\typeof{N_1}{C}}
        {\typeof{x_2}{A_2}}{\typeof{N_2}{C}},
    \]
    असेल, तर $M' \ident \Subst{N_i}{N}{\typeof{x_i}{A_i}}$.
    ~$M'$ मधील रेडेक्स घ्या. एकतर तो बसवलेल्या~$N$ च्या प्रतीत आहे; अशा वेळी
    गृहीतकामुळे त्याची कट-कोटी मूळ~$M$ पेक्षा कमी आहे.
    अन्यथा~$N_i$ मध्ये त्याला
    अनुरूप व समान कट-कोटीचा रेडेक्स आहे; गृहीतकामुळे ती
    कट-कोटी $\cutrank{M}$ पेक्षा कमी आहे. अन्यथा त्याची
    कट-कोटी $\len{A_i}$ आहे; व्याख्येनुसार ती
    $\cutrank{\ore{\inj[A_{3-i}]{i}{\typeof{N}{A_i}}}
      {\typeof{x_1}{A_1}}{\typeof{N_1}{C}}
      {\typeof{x_2}{A_2}}{\typeof{N_2}{C}}}$ पेक्षा कमी आहे.
  \end{enumerate}
\end{proof}

\begin{thm}
  प्रत्येक योग्य प्रकाराचे सिद्धता-पद पूर्ण सामान्यरूपीकरणाच्या
  अनुरूप पद-रूपांतरणांनी सामान्य रूपात न्यूनीत करता येते;
  प्रत्येक अंतःप्रज्ञावादी निष्पत्ती सामान्य निष्पत्तीमध्ये न्यूनीत करता येते.
\end{thm}

\begin{proof}
  नैसर्गिक निगमनाच्या सामान्यरूपीकरण प्रकरणातील
  \ref{pt:nor:thm:thm:normalization-N2i} वापरू.
  योग्य प्रकाराच्या~$M$ साठी त्याची प्रकार-सिद्धता घ्या.
  अनावश्यक संदर्भ-गृहीतके काढून आणि बद्ध चलांची नावे
  सुरक्षित ठेवून तिच्याशी अनुरूप \Log{N2i}-सिद्धता
  मिळते. तिचे सामान्यरूप करा.

  प्रत्येक रूपांतरणाला पुन्हा सिद्धता-पदे द्या. मुख्य
  रूपांतरणे म्हणजे येथे दिलेले फलन, गुणाकार आणि बेरीज
  प्रकारांचे बीटा-नियम. वियोग-निरसनापुढील निरसन आत
  हलवणे हे संबंधित क्रमपरिवर्तन आहे; त्यात बद्ध चल
  आत हलवलेल्या पदांच्या मुक्त चलांपासून नवे ठेवायचे.
  \FalseInt{} च्या सरलीकरणासाठीही त्या प्रकरणातील
  रूपांतरणाची अनुरूप पद-आवृत्ती घ्यायची. यामुळे मूळ
  पदापासून सामान्य सिद्धतेच्या पदापर्यंत सांत क्रम मिळतो.
  अंतिम पदात लागू होणारे रूपांतरण असते, तर अंतिम
  सिद्धता सामान्य नसती. आवश्यकतेनुसार उरलेला संदर्भ
  पुन्हा वाढवून \Log{tN3ip} मध्ये मूळ प्रकार जपता येतो.

  म्हणून दोन्ही दावे मिळतात. वरचा रेडेक्स-मापदंड
  एकटाच संपूर्ण सिद्धता देत नाही: एका उपपदाचे संकुचन
  बाहेरच्या संदर्भात नवा रेडेक्स उघडू शकते; विशेषतः
  वियोग-निरसनाच्या निष्कर्षाचा प्रकार त्याच्या वियोगाच्या
  प्रकारापेक्षा मोठा असू शकतो. म्हणून नुसती महत्तम
  कट-कोटीच्या रेडेक्सांची संख्या घटते असे गृहीत धरलेले नाही.
\end{proof}

प्रकारयुक्त $\lambd$-पदांचे सामान्यीकरण होते याचा अर्थ
अशा प्रत्येक पदाला \emph{सामान्य रूप मिळते}. $\lambd$-पदांचे
सामान्य रूपापर्यंत $\beta$-न्यूनीकरण करणे म्हणजे संगणन चालवणे,
आणि सामान्य रूप म्हणजे त्याचे मूल्य, असे मानल्यास वरच्या पूर्ण सामान्यरूपीकरणाच्या पद्धतीने प्रत्येक
योग्य प्रकाराचे पद अखेरीस सामान्य रूपात पोहोचते.
शिवाय प्रकार-संरक्षणामुळे त्या मूल्याचा प्रकार योग्य असतो.
उदाहरणार्थ,~$N$ चा प्रकार $A \lif B$ असेल
(म्हणजे ते प्रकार~$A$ च्या आर्ग्युमेंटवर चालणारे आणि
प्रकार~$B$ चे मूल्य देणारे फलन असेल) आणि प्रकार~$A$
च्या~$M$ या पदावर (निवेशावर) ते लावले, तर $(NM)$ चे
$N'$ या पदात (निर्गमात) न्यूनीकरण होते आणि त्या निर्गमाचा
प्रकार~$B$ असेल याची हमी मिळते.

मुख्य बीटा-न्यूनीकरणासाठी प्रकारयुक्त $\lambd$-कलनाचा
आणखी एक गुणधर्म खरा आहे: बीटा-पायऱ्या कोणत्याही क्रमाने
लावल्या तरी योग्य प्रकाराच्या पदातून अनंत क्रम मिळत नाही.
याला \emph{प्रबळ सामान्यीकरण} म्हणतात. हा अधिक बलवान
परिणाम येथे स्वतंत्रपणे सिद्ध केलेला नाही; वरचे प्रमेय
एका पूर्ण सामान्यरूपीकरण-क्रमाचे अस्तित्व देते.
खालील संगमशीलतेचा विचार मुख्य बीटा-न्यूनीकरणाचा आहे;
सर्व वाढीव क्रमपरिवर्तनांचे स्वतंत्र संगमशीलता-प्रमाणपत्र नाही.

संगणन चालवण्याच्या वेगवेगळ्या पद्धतींनी \emph{तेच}
मूल्य मिळते हे कसे कळते? आतापर्यंत (प्रकार-संरक्षणामुळे)
एवढेच माहीत आहे की मिळालेल्या सर्व मूल्यांचा प्रकार एकच
असतो. प्रकारयुक्त $\lambd$-कलनाला आणखी एक महत्त्वाचा
गुणधर्म आहे: ते \emph{संगमशील}, म्हणजे त्याला
\emph{चर्च–रॉसर गुणधर्म} आहे. $N \red N_1$ आणि
$N \red N_2$ असतील, तर एखादे पद~$N'$ असे असते की
$N_1 \red N'$ आणि $N_2 \red N'$. लक्षात ठेवा की
$N \red N'$ याचा अर्थ $\redone$-न्यूनीकरणाच्या
शून्य किंवा अधिक पायऱ्यांनी~$N'$ मिळते.
$N_1$ सामान्य रूपात असेल, तर $N_1 \red N'$ असणारे
एकमेव पद~$N'$ म्हणजे~$N_1$ स्वतःच (शून्य
$\redone$-न्यूनीकरण पायऱ्यांनी). त्यामुळे~$N_1$ आणि
$N_2$ दोन्ही सामान्य रूपात असतील, तर $N_1 = N' = N_2$.
दुसऱ्या शब्दांत, $\redone$ प्रबळ सामान्यीकारी असल्यामुळे
पद न्यूनीत करण्याची प्रत्येक पद्धत अखेरीस सामान्य रूप देते.
$\red$ ला चर्च–रॉसर गुणधर्म असल्यामुळे कोणतीही दोन
सामान्य रूपे समान असतात. म्हणून प्रकारयुक्त
$\lambd$-कलनातील संगणन नेहमी थांबते आणि ते कसेही
चालवले तरी मिळणारे मूल्य (सामान्य रूप) तेच असते.

संबंध संगमशील आहे हे सिद्ध करणे सहसा अवघड असते.
तथापि, पुढील दुर्बल अट सहज सिद्ध करता येते:

\begin{prop}[दुर्बल चर्च–रॉसर गुणधर्म]
$N \redone N_1$ आणि $N \redone N_2$ असतील, तर एखादे
पद~$N'$ असे असते की $N_1 \red N'$ आणि $N_2 \red N'$.
\end{prop}

\begin{proof}
  $N$ मधील रेडेक्स~$M_1$ किंवा~$M_2$ त्याच्या संकुचनफलाने
  बदलल्यास अनुक्रमे~$N_1$ किंवा~$N_2$ मिळतात.
  $M_1$ आणि~$M_2$ हे~$N$ चे उपपद असल्याने त्यांचे
  अंशतः छेदन होऊ शकत नाही. म्हणजे एक रेडेक्स दुसऱ्याचे
  उपपद असेल, किंवा ते~$N$ मधील वेगवेगळ्या न छेदणाऱ्या
  स्थानी असतील. दुसरी बाब घ्या:~$N$ चे रूप $N(M_1,M_2)$
  आहे. $M_1$ च्या जागी त्याचे संकुचनफल~$M_1'$ ठेवल्यास
  $N(M_1',M_2)$ हे~$N_1$ मिळते; $M_2$ च्या जागी त्याचे
  संकुचनफल~$M_2'$ ठेवल्यास $N(M_1,M_2')$ हे~$N_2$
  मिळते. दोघांचेही $N(M_1',M_2')$ मध्ये न्यूनीकरण होते.
  दुसऱ्या बाबतीत~$M_2$ हे~$M_1$ चे उपपद आहे असे माना
  (उलट बाब सममित आहे), आणि~$M_1$ कोणत्या प्रकारचा
  रेडेक्स आहे यानुसार बाबी वेगळ्या करा. उदाहरणार्थ,
  $M_1 \ident \proj{1}{\pair{O_1}{O_2}}$ असेल, तर त्याचे
  $O_1$ मध्ये न्यूनीकरण होते. $M_2$ हे~$O_1$ चे उपपद
  असेल, तर~$M_2$ चे परिवर्तन केल्यावर~$O_1'$ मिळते.
  म्हणून आधी~$M_1$ आणि नंतर~$M_2$ चे परिवर्तन केल्यावर
  $O_1'$ मिळते. पण आधी~$M_2$ चे परिवर्तन केल्यावर
  $\proj{1}{\pair{O_1'}{O_2}}$ मिळते, आणि त्याचेही
  $O_1'$ मध्ये न्यूनीकरण होते.
\end{proof}
दोन्ही निवडी तोच रेडेक्स असतील, तर निष्कर्ष तत्काळ
मिळतो. वरील गुणाकार उदाहरणात~$M_2$ हे~$O_2$ मध्ये
असेल, तर दोन्ही क्रम~$O_1$ येथे मिळतात. फलनाच्या
बीटा-रेडेक्समध्ये आतील पायरी देहात असेल, तर
प्रतिस्थापनानंतर ती अनुरूप स्थानी लावायची; ती आदानात
असेल, तर बसवलेल्या सर्व प्रतींत लावायची, किंवा आदान
वापरले नसेल, तर शून्य पायऱ्या घ्यायच्या. बेरीज-प्रकारात
निवडलेल्या शाखेसाठी हाच युक्तिवाद लागू होतो; टाकलेल्या
शाखेतील पायरीची गरज उरत नाही. सुरक्षित प्रतिस्थापन
आणि बद्ध चलांची सुसंगत नावे वापरल्याने सामायिक पद
मिळते. पदांची समानता येथे बद्ध-चल नावबदलापर्यंत आहे.

\begin{lem}[न्यूमनची प्रमेयिका]
$\redone$ प्रबळ सामान्यीकारी असेल आणि त्याला दुर्बल
चर्च–रॉसर गुणधर्म असेल, तर त्याचा शून्य किंवा अधिक पायऱ्यांचा संबंध~$\red$ संगमशील आहे.
\end{lem}

\begin{proof}
  $\redone$ प्रबळ सामान्यीकारी असल्यामुळे प्रत्येक पूर्ण
  न्यूनीकरण-क्रम सामान्य रूपात संपतो. सामान्य रूपे एकमेव असणे आणि
  $\red$ संगमशील असणे या सममूल्य अटी आहेत. म्हणून एखाद्या
  पद~$N$ ला~$N'$ आणि~$N''$ ही दोन वेगळी सामान्य रूपे
  आहेत असे माना; ती पुढील न्यूनीकरण-क्रमांनी मिळतात:
  \begin{align*}
  N & \redone N_1' \redone N_2' \redone \dots \redone N_k' = N' \text{ आणि}\\
  N & \redone N_1'' \redone N_2'' \redone \dots \redone N_l'' = N''
  \end{align*}
  (न्यूनीकरण-क्रमाची लांबी $>0$ असलीच पाहिजे, अन्यथा
  $N$ आधीच सामान्य रूपात असेल.)

  $N_1' = N_1''$ असेल, तर त्याला दोन वेगळी सामान्य रूपे
  ($N'$ आणि $N''$) आहेत. $N_1' \neq N_1''$ असेल, तर
  दुर्बल चर्च–रॉसर गुणधर्मामुळे एखादे~$N'''$ असे आहे की
  $N_1' \red N'''$ आणि $N_1'' \red N'''$.
  $N' \neq N''$ असल्यामुळे $N''' \neq N'$ किंवा
  $N''' \neq N''$ यांपैकी किमान एक अट खरी आहे.
  निर्णायक मुद्दा असा: प्रबळ सामान्यीकरणामुळे या सामायिक
  पदाला एक सामान्य रूप मिळते. ते मूळ दोन वेगळ्या सामान्य
  रूपांपैकी किमान एकापेक्षा वेगळे असते. त्यामुळे~$N_1'$
  किंवा~$N_1''$ यांपैकी एकाला दोन वेगळी
  सामान्य रूपे आहेत. आपण दाखवले की~$N$ ला दोन वेगळी
  सामान्य रूपे असतील, तर एखादे~$N_1$ असे आहे की
  $N \redone N_1$ आणि~$N_1$ लाही दोन वेगळी सामान्य रूपे
  आहेत. ही प्रक्रिया पुढे चालू ठेवून
  $N \redone N_1 \redone N_2 \redone \dots$ असा क्रम
  मिळेल; पण $\redone$ प्रबळ सामान्यीकारी असल्यामुळे
  हे अशक्य आहे.
\end{proof}



